%% ma: L10-B15-0093
%% lop: 10
%% chuong: 6
%% bai: 15
%% dang: 10.15.4
%% loai: TN4
%% muc: NB
%% dap_an: "B"
%% nguon: "Ảnh giáo viên gửi 25/11/2026 (lô 2): Bài tập trắc nghiệm – Toán 10 Bài 15 Hàm số (sự biến thiên, đồ thị), Phần I Câu 4"
%% kien_thuc: "Lớp 10 Bài 15 (đồ thị, sự đồng biến, nghịch biến của hàm số)"
%% trang_thai: chinh-thuc
%% ghi_chu: "Đồ thị vẽ lại theo hàm y = -x^4 + 2x^2 + 1 (khớp các điểm (±1;2), (0;1) của hình gốc)."
\begin{ex}
Cho hàm số $y=f(x)$ có đồ thị là đường cong trong hình bên.
\begin{center}\begin{tikzpicture}[scale=.9]
\draw[truc] (-2,0)--(2.2,0) node[below]{$x$};
\draw[truc] (0,-.5)--(0,2.8) node[left]{$y$};
\draw[phu] (-1,0)--(-1,2)--(1,2)--(1,0) (0,1)--(.15,1);
\node[below] at (-1,0) {\footnotesize$-1$}; \node[below] at (1,0) {\footnotesize$1$};
\node[below left] at (0,0) {\footnotesize$O$};
\node[above right] at (0,2) {\footnotesize$2$}; \node[below right] at (0,1) {\footnotesize$1$};
\draw[thick,domain=-1.55:1.55,smooth,samples=80] plot (\x,{-(\x)^4+2*(\x)^2+1});
\end{tikzpicture}\end{center}
Hàm số đã cho đồng biến trên khoảng nào dưới đây?
\choice{$(-\infty;0)$}{\True $(0;1)$}{$(1;+\infty)$}{$(-1;0)$}
\loigiai{Trên $(0;1)$ đồ thị đi lên nên hàm số đồng biến trên $(0;1)$.}
\end{ex}
